\section{智能驾驶 vs 具身智能：视觉游戏与物理交互}

上一章讲了合成数据如何在自动驾驶中起作用，本章转向具身智能为什么更难。自动驾驶不是简单问题，但它有一个巨大优势：车这个平台已经存在上百年，大量车辆在真实道路上运行，L2+ 车队能持续回传真实视觉、传感器和驾驶数据。机器人则没有上百万台人形或四足机器人在家庭、工厂、酒店里自然运行，因此真实数据飞轮很难先转起来。

这也是谢晨判断“具身智能必须首先依赖大量合成数据，再用少量真实数据校准”的原因。自动驾驶更像视觉和规则游戏：道路、车道、交通参与者、车辆本体相对明确；具身智能要处理杯子、吸管、冰箱、键盘、布料、液体、摩擦、碰撞和不同机器人本体。物理交互的复杂度，使它不能只靠互联网上已有的图像和视频。

\begin{figure}[H]
\centering
\includegraphics[width=0.94\textwidth]{av-vs-embodied-data.png}
\caption{智能驾驶 vs 具身智能：智能驾驶偏视觉游戏，具身智能更依赖物理交互。自制概念图，依据 00:16:17--00:32:41 对谈内容整理。}
\end{figure}

\begin{knowledgebox}{读图：为什么自动驾驶的真实数据飞轮更容易}
左侧自动驾驶有大量真实车辆、稳定本体和道路结构，因此真实数据占主体，合成数据主要放大长尾和安全问题。右侧具身智能缺少自然运行的机器人群体，任务又跨物体、空间、接触和动作，因此合成数据与仿真必须提前承担“造世界”的角色。
\end{knowledgebox}

\subsection{物理交互瓶颈}

本节进一步解释“物理交互”到底增加了什么难度。视觉数据在互联网上非常丰富，照片和视频可以帮助生成街道、车辆、树木和路面；但物理交互数据不只是外观，而是力、摩擦、转轴、碰撞、材质、重量、形变和失败恢复。没有底层交互数据，AI 很难凭空生成可靠的物理行为。

冰箱门是访谈里的好例子。要让机器人学会开冰箱，仿真资产不能只是一个冰箱外壳。它需要门的转轴、铰链、磁吸力、不同角度下的开门力、抽屉和门板的碰撞体、不同冰箱型号的分布，以及机器人手和门之间的接触力。视觉上像冰箱不够，物理上可交互才是关键。

\begin{figure}[H]
\centering
\includegraphics[width=0.96\textwidth]{physical-interaction-bottleneck.png}
\caption{物理交互瓶颈：抓取、接触、力反馈和材料差异让数据特别贵。自制概念图，依据 00:16:17--00:32:41 对谈内容整理。}
\end{figure}

\begin{knowledgebox}{读图：具身数据的单位不是 frame，而是 interaction}
图中每一个变量都能导致任务失败：抓取点稍偏会滑落，材质不同会改变摩擦，执行器误差会累积，碰撞体不准会让仿真策略在真实世界失效。对具身智能来说，一条数据如果没有动作、力学结果和失败反馈，就很难教会机器人真正做事。
\end{knowledgebox}

\begin{warningbox}{不要把视频数据误当作机器人数据}
互联网视频能告诉模型“人类怎么做”，但通常没有机器人本体、关节状态、接触力、动作命令和环境可控变量。它适合预训练世界知识，却不能直接替代机器人训练所需的闭环交互数据。
\end{warningbox}

\subsection{遥操作路线与商业闭环}

本节回应一个常见反问：既然真实数据重要，能不能通过遥操作机器人先赚钱、再采数据？谢晨认为这在少数场景可能成立，但作为全局路线很难快过合成数据为主的路线。原因是遥操作必须先产生客户愿意付费的价值，否则只是高成本采样；跨州、跨国遥操作还会遇到监管、人力成本、工会、场景租赁和运营复杂度。

自动驾驶的数据飞轮之所以强，是因为车本来就能卖，用户开车时顺手把数据带回来；机器人还没有进入这个状态。一个遥操作机器人如果不能在机场、酒店、餐饮、零售等场景创造足够 ROI，它就不能形成“付费使用产生数据，数据提升能力，能力带来更多使用”的正循环。讲者判断，特斯拉或许有内部主机厂场景和组织能力走一部分路线，但多数机器人公司不能简单复制。

\begin{knowledgebox}{术语消化：真实路线、遥操作与合成路线}
\begin{tabularx}{\textwidth}{p{0.20\textwidth}p{0.35\textwidth}X}
\toprule
路线 & 数据来源 & 主要约束 \\
\midrule
真实车队路线 & 量产设备自然运行并回传 & 需要产品先有用，硬件规模已经存在。 \\
遥操作路线 & 人远程控制机器人并采集动作轨迹 & 人力、场景、法规和 ROI 压力很高。 \\
合成数据路线 & 仿真场景中批量生成任务和交互 & 需要高质量 Real2Sim、物理引擎和评价系统。 \\
混合路线 & 真实少量校准，合成大规模扩增 & 当前最可能成为主流的工程方向。 \\
\bottomrule
\end{tabularx}
\end{knowledgebox}

\subsection{本章小结}

自动驾驶的数据瓶颈主要集中在长尾和安全场景；具身智能的数据瓶颈则深入到物理交互本身。机器人没有现成的大规模真实数据飞轮，因此合成数据和物理仿真不是补丁，而是早期路线的核心。

